\documentclass[10pt,a4paper]{amsart}
\usepackage[margin=19mm]{geometry}
\usepackage{amsmath,amssymb,mathtools}
\usepackage[scr=rsfs,cal=euler]{mathalfa}
\usepackage{xfrac}
\usepackage[unicode,hidelinks,bookmarks=false]{hyperref}
\newcommand{\ie}{i.e.\ }
\newcommand{\cf}{cf.\ }
\newcommand{\resp}{respectively\ }
\newcommand{\Subsection}{\subsection}
\providecommand{\nobreakdash}{\nobreak\hbox{-}}
\setlength{\emergencystretch}{2em}
\multlinegap0pt
\usepackage{lmodern}
\usepackage{mathtools}
\usepackage{mathrsfs}
\usepackage{enumitem}
\usepackage{booktabs}
\usepackage{tabularx}
\usepackage{longtable}
\usepackage{cleveref}
\allowdisplaybreaks
\numberwithin{equation}{section}
\renewcommand{\arraystretch}{1.14}
\newtheorem{theorem}{Theorem}[section]
\newtheorem{proposition}[theorem]{Proposition}
\newtheorem{lemma}[theorem]{Lemma}
\newtheorem{corollary}[theorem]{Corollary}
\newtheorem{definition}[theorem]{Definition}
\newtheorem{assumption}[theorem]{Assumption}
\newtheorem{conjecture}[theorem]{Conjecture}
\newtheorem{problem}[theorem]{Problem}
\newtheorem{principle}[theorem]{Principle}
\newtheorem{remark}[theorem]{Remark}
\DeclareMathOperator*{\esssup}{ess\,sup}
\DeclareMathOperator{\dist}{dist}
\DeclareMathOperator{\supp}{supp}
\DeclareMathOperator{\Range}{Range}
\DeclareMathOperator{\Ker}{Ker}
\DeclareMathOperator{\Span}{span}
\DeclareMathOperator{\Tr}{tr}
\newcommand{\R}{\mathbb R}
\newcommand{\N}{\mathbb N}
\newcommand{\eps}{\varepsilon}
\newcommand{\CKN}{\mathrm{CKN}}
\newcommand{\NS}{\mathrm{NS}}
\newcommand{\loc}{\mathrm{loc}}
\newcommand{\dx}{\,dx}
\newcommand{\dt}{\,dt}
\newcommand{\dxdt}{\,dx\,dt}
\newcommand{\norm}[2]{\left\|#1\right\|_{#2}}
\newcommand{\abs}[1]{\left|#1\right|}
\newcommand{\ip}[2]{\left\langle #1,#2\right\rangle}
\newcommand{\Sup}{\mathrm{Sup}}
\newcommand{\Tax}{\mathrm{Tax}}
\newcommand{\Leak}{\mathrm{Leak}}
\newcommand{\Rep}{\mathrm{Rep}}
\newcommand{\Err}{\mathrm{Err}}
\newcommand{\Obs}{\mathrm{Obs}}
\newcommand{\calA}{\mathcal A}
\newcommand{\calB}{\mathcal B}
\newcommand{\calC}{\mathcal C}
\newcommand{\calD}{\mathcal D}
\newcommand{\calE}{\mathcal E}
\newcommand{\calF}{\mathcal F}
\newcommand{\calG}{\mathcal G}
\newcommand{\calH}{\mathcal H}
\newcommand{\calK}{\mathcal K}
\newcommand{\calL}{\mathcal L}
\newcommand{\calM}{\mathcal M}
\newcommand{\calN}{\mathcal N}
\newcommand{\calO}{\mathcal O}
\newcommand{\calP}{\mathcal P}
\newcommand{\calQ}{\mathcal Q}
\newcommand{\calR}{\mathcal R}
\newcommand{\calS}{\mathcal S}
\newcommand{\calT}{\mathcal T}
\newcommand{\calX}{\mathcal X}
\newcommand{\calY}{\mathcal Y}
\newcommand{\calZ}{\mathcal Z}
\begin{document}
\title[A Structural Audit of Navier-Stokes Obstruction Calculus]{A Structural Audit of Navier-Stokes Obstruction Calculus}
\author{Runlong Yu}
\date{}
\maketitle
\noindent\emph{Source and licence.} A Structural Audit of Navier-Stokes Obstruction Calculus. Version arXiv:2606.25341v1 (CC BY 4.0). This material is distributed under the Creative Commons Attribution 4.0 International licence, http://creativecommons.org/licenses/by/4.0/, as recorded in the arXiv OAI licence metadata for this exact version and shown on the abstract page https://arxiv.org/abs/2606.25341v1. Full rights notice: licenses/arxiv-2606.25341-CC-BY-4.0.txt.\par\smallskip
\noindent\textit{Editorial scope.} Counted unit: the original \texttt{theorem} environment \texttt{thm:finite-gap} at source lines 523--548 together with its complete proof at lines 550--552; the delivered document keeps source lines 501--553 so that every referenced label and every citation resolves inside it. Native editable source is the paper's own concatenated TeX; the AMS wrapper is editorial formatting only. No publisher class, upstream script or unrelated artwork is redistributed.
\section{Finite-window obstruction geometry}\label{sec:finite-window}

\subsection{Cleaning, quotient, and observation}

Fix a finite scale or frequency window $\Lambda$. Let $D_\Lambda$ be a finite-dimensional constrained defect space, $C_\Lambda$ a cleaning space, and
\[
G_\Lambda:C_\Lambda\to D_\Lambda
\]
a linear cleaning map. The quotient distance is
\begin{equation}\label{eq:quot-distance}
\dist_\Lambda(d,\operatorname{Im}G_\Lambda)
=
\inf_{c\in C_\Lambda}\norm{d-G_\Lambda c}{D_\Lambda}.
\end{equation}
Let
\[
O_\Lambda:D_\Lambda\to Y_\Lambda
\]
collect active pressure, flux, energy, and trace observations. In a dynamic version, one also has a reproduction residual $\Rep_\Lambda(d)$ and a nonnegative tax or profit-cost functional $\Tax_\Lambda(d)$.

\subsection{Exact finite-window quotient theorem}


\begin{theorem}[Finite-window anti-phantom gap]\label{thm:finite-gap}
Assume the detector
\begin{equation}\label{eq:detector-size}
\calM_\Lambda(d)
=\norm{O_\Lambda d}{Y_\Lambda}
+C_R\Rep_\Lambda(d)+C_T\Tax_\Lambda(d)
\end{equation}
is continuous, positively homogeneous, and constant on gauge cosets. Assume further that
\begin{equation}\label{eq:kernel-free}
O_\Lambda d=0,
\quad
\Rep_\Lambda(d)=0,
\quad
\Tax_\Lambda(d)=0
\quad\Longrightarrow\quad
d\in\operatorname{Im}G_\Lambda.
\end{equation}
Then there exists $\mu_\Lambda>0$ such that
\begin{equation}\label{eq:finite-gap}
\calM_\Lambda(d)
\ge
\mu_\Lambda\dist_\Lambda(d,\operatorname{Im}G_\Lambda)
\qquad
\text{for all }d\in D_\Lambda.
\end{equation}
\end{theorem}

\begin{proof}
The detector descends to a continuous positively homogeneous function on the finite-dimensional quotient $D_\Lambda/\operatorname{Im}G_\Lambda$. By \eqref{eq:kernel-free}, it is strictly positive on the quotient unit sphere. Compactness of that sphere gives a positive minimum $\mu_\Lambda$. This is a finite-dimensional compactness step; it should not be confused with scale-uniform PDE compactness such as the Aubin--Lions--Simon mechanism \cite{Simon1986}. Homogeneity gives \eqref{eq:finite-gap}.
\end{proof}


\begin{thebibliography}{99}
\bibitem[Simon1986]{Simon1986} J.~Simon, \newblock Compact sets in the space $L^p(0,T;B)$, \newblock \emph{Annali di Matematica Pura ed Applicata} \textbf{146} (1986), 65--96. \newblock DOI: \url{https://doi.org/10.1007/BF01762360}.
\end{thebibliography}
\end{document}
